% Generated by tools/edn_latex.py from reports/edn.md.
% Do not edit: change reports/edn.md and run `make edn`.
\ednentry{
  id = {EDN-60},
  title = {Which segments may gate is enforced in code that fails, not by a rule in a document},
  date = {2026-09-30},
  milestone = {M3: Quality Assurance},
  topic = {Model Evaluation (SC-04), Testing Strategy},
  participants = {@lukas2510},
  decision = {The three rules that decide which segment levels SC-04 is allowed to gate on are checks that fail the stage rather than conventions a reader has to remember. A module constant names the segments that condition on the target (\texttt{price\_\allowbreak{}bucket}), and adding one to \texttt{evaluate.\allowbreak{}segments} raises; the row builder reads that constant itself, so no caller can produce a price-bucket row that counts. A segment over an input field FR-01 does not require also raises, because the rule that excludes a \texttt{\textquotedbl{}(missing)\textquotedbl{}} level is only correct for a required field. And a per-make level naming a make the model was not fitted on raises, so EDN-48's third obligation is a checked property rather than an intention. Every excluded level stays in the exported table with its metrics, its row count and the name of the rule that excluded it.},
  rationale = {\emph{Alternatives considered:}
\begin{itemize}[leftmargin=*, topsep=2pt]
\item \textbf{Option A (chosen): a constant plus three refusals, and every exclusion recorded in the table.}
\newline Pros: the rule lives where it is applied, so it cannot be forgotten by someone editing \texttt{params.\allowbreak{}yaml} a month from now, and each refusal names the reason rather than the symptom. Keeping the excluded rows in the table with \texttt{excluded\_\allowbreak{}because} means the report can show that \texttt{Volkswagen} had 81 test rows and \texttt{fuel\_\allowbreak{}category=\allowbreak{}(missing)} had 5, which is the difference between a breakdown a reader can audit and one that silently omits what did not fit.
\newline Cons: more code than a comment, and three error paths that only fire on a configuration nobody has written yet.
\item \textbf{Option B: document the rules in the problem specification and rely on review.}
\newline Pros: no code, and the specification is where the criterion is defined anyway.
\newline Cons: the failure it has to prevent is a one-line \texttt{params.\allowbreak{}yaml} edit that produces a passing gate. SC-04's worst qualifying level for \texttt{lgbm-\allowbreak{}extended} is 17.8\,\%; its worst price bucket is 22.3\,\%, so adding \texttt{price\_\allowbreak{}bucket} to the segment list changes the number the gate blocks on while looking like a reporting improvement. A rule that is only written down is checked exactly as often as somebody reads it.
\item \textbf{Option C: drop the price buckets from the stage entirely.}
\newline Pros: nothing to exclude, so nothing to enforce.
\newline Cons: they are the diagnostic that shows \emph{why} they may not gate. Measured on the real snapshot, their profile is U-shaped (22.3\,\% under 5,000 EUR and 6.3\,\% over 80,000 EUR against 5.5\,\% in the 40,000 to 80,000 EUR bucket), which is the regression-to-the-mean artefact of conditioning on the target. Without the table the exclusion is an assertion; with it, it is evidence.
\item \textbf{Option D: exclude a \texttt{\textquotedbl{}(missing)\textquotedbl{}} level for every segment, without checking that the field is required.}
\newline Pros: one rule, no second check.
\newline Cons: it would be wrong the first time somebody segments by an optional field. For a required field the absence cannot occur at serving time, because FR-01 answers such a request with a 422; for an optional field it is a case the API answers every day, and excluding it would hide exactly the rows worth looking at.
\end{itemize}
\par
\emph{Why:} All three rules share one shape: a statement about what the criterion means that would otherwise be maintained by discipline. The M3 rubric rewards exactly this distinction, and the measurement makes it concrete rather than stylistic: the price buckets would change SC-04's verdict if they gated, so the difference between option A and option B is a number in the report. Option D was rejected on the same grounds as the price buckets: it is right today only because all five configured segments happen to be required fields, and an accidental correctness is what the check turns into a guaranteed one.},
  ai = {Alternative generation, Alternative assessment, Recommendation, Solution generation},
  response = {Accepted with modifications},
  assessment = {AI proposed the target-conditioned constant and the ``absence of a required field'' exclusion, and measured the price-bucket profile that shows why the exclusion matters rather than asserting that it does. Two modifications were made to what it proposed. Its construction passed an \texttt{eligible=\allowbreak{}False} flag into the row builder for the price-bucket rows, which leaves the exclusion dependent on the caller passing the flag; the flag was removed and the builder now reads the segment's own name, so the exclusion holds however the rows are built. And it did not propose the third check at all: the requirement that a segment be a required input field came from asking what makes the \texttt{\textquotedbl{}(missing)\textquotedbl{}} exclusion correct, which turned out to be an assumption nothing enforced.},
  aievidence = {Claude Code session on 2026-09-30 designing issue \#39, which measured the price-bucket breakdown at 22,551 test rows and reported ``the U shape is the regression-to-the-mean artefact of conditioning on the target, not a fairness finding, and the 490k MAE in the top bucket would dominate any aggregate''; and the implementation session on the same day, which replaced the \texttt{eligible} flag with the name-based rule while reducing the row builder's argument count.},
  evidence = {\href{https://github.com/mlops-2627q1-mds-upc/MLOPS_Recommenditos/blob/main/recommenditos/modeling/evaluate.py}{\texttt{recommenditos/\allowbreak{}modeling/\allowbreak{}evaluate.\allowbreak{}py}}, \texttt{TARGET\_\allowbreak{}CONDITIONED\_\allowbreak{}SEGMENTS}, \texttt{criterion\_\allowbreak{}segments}, \texttt{\_\allowbreak{}sc04\_\allowbreak{}eligibility} and \texttt{check\_\allowbreak{}make\_\allowbreak{}levels\_\allowbreak{}are\_\allowbreak{}served}; \texttt{tests/\allowbreak{}test\_\allowbreak{}model.\allowbreak{}py::test\_\allowbreak{}a\_\allowbreak{}target\_\allowbreak{}conditioned\_\allowbreak{}segment\_\allowbreak{}cannot\_\allowbreak{}become\_\allowbreak{}a\_\allowbreak{}criterion}, \texttt{::test\_\allowbreak{}a\_\allowbreak{}segment\_\allowbreak{}over\_\allowbreak{}an\_\allowbreak{}optional\_\allowbreak{}input\_\allowbreak{}field\_\allowbreak{}is\_\allowbreak{}refused}, \texttt{::test\_\allowbreak{}a\_\allowbreak{}price\_\allowbreak{}bucket\_\allowbreak{}row\_\allowbreak{}is\_\allowbreak{}reported\_\allowbreak{}and\_\allowbreak{}cannot\_\allowbreak{}count}, \texttt{::test\_\allowbreak{}sc04\_\allowbreak{}refuses\_\allowbreak{}rows\_\allowbreak{}from\_\allowbreak{}a\_\allowbreak{}segment\_\allowbreak{}it\_\allowbreak{}may\_\allowbreak{}not\_\allowbreak{}gate}, \texttt{::test\_\allowbreak{}a\_\allowbreak{}per\_\allowbreak{}make\_\allowbreak{}level\_\allowbreak{}the\_\allowbreak{}api\_\allowbreak{}would\_\allowbreak{}reject\_\allowbreak{}is\_\allowbreak{}refused}; \texttt{reports/\allowbreak{}metrics/\allowbreak{}segments.\allowbreak{}csv}; \href{https://github.com/mlops-2627q1-mds-upc/MLOPS_Recommenditos/blob/main/docs/docs/model-card.md}{model card}, the per-segment breakdown; \hyperref[EDN-48]{EDN-48}; \href{https://github.com/mlops-2627q1-mds-upc/MLOPS_Recommenditos/issues/39}{issue \#39}.},
}
